
\chapter{Диспуты и чарджбэки}
\label{ch:disputes}
\margintoc


\begin{chapteroverview}
  Чарджбэк встроен в карточную систему намеренно: он даёт держателю карты
  право оспорить уже завершённую операцию и переводит спор в формальный
  цикл доказывания. В этой главе мы разберём путь спора от первой жалобы
  до арбитража, увидим, как \scheme{Visa} и \scheme{Mastercard}
  по-разному называют сходные стадии, поймём логику кодов основания
  спора и требований к доказательствам, а затем отдельно посмотрим,
  зачем \scheme{Visa} ввела Compelling Evidence 3.0 и почему у СБП спор
  устроен иначе.
\end{chapteroverview}

Когда держатель карты звонит в банк и говорит: \enquote{я этого не покупал},
вчерашняя завершённая транзакция перестаёт быть спокойным прошлым. Она
превращается в спор, где на каждой стадии заново распределяются деньги,
сроки и бремя доказывания.

\section{Жизнь диспута: пять стадий}
\label{sec:disputes-lifecycle}

Что на самом деле делает банк, когда получает жалобу? Он не принимает
мгновенное решение о возврате. Он запускает регламентированный цикл, в
котором каждая стадия задаёт свой срок, свой набор документов и свои
финансовые последствия
\sidecite{visa-dispute-rules}\sidecite{mc-chargeback-guide}.

Сначала нередко возникает предварительный запрос на сведения. В отрасли
его называют inquiry или retrieval request: эмитент или его процессор
пытается понять, есть ли основания переводить ситуацию в формальный
спор. Этот этап полезен, потому что часть недоразумений снимается ещё
до чарджбэка, но обязательным он не является.

Если спор не снят, начинается формальный чарджбэк (chargeback):
эмитент обращается к схеме, схема дебетует эквайера, а эквайер уже
переносит риск на мерчанта. К спору прикладывается код основания
спора (reason code): именно он показывает, что нужно доказывать и на
чьей стороне лежит основная обязанность по доказательствам
\sidecite{visa-dispute-rules}.

Далее следует ответ мерчанта на чарджбэк, который в отраслевой речи
зовётся ответом эквайера на чарджбэк (representment). В \scheme{Visa}
официальная терминология ближе к словарю \enquote{диспут / ответ на
диспут / предарбитраж / арбитраж}, а в \scheme{Mastercard}
по-прежнему заметна логика циклов чарджбэка и подачи дела. Смысл,
однако, одинаков: мерчант должен приносить доказательства, а не
просто возражения.

Если ответ не убедил эмитента, наступает предарбитраж (pre-arbitration):
эмитент либо принимает ответ, либо поднимает спор ещё раз, уже на более
узкой и дорогой стадии. Если и это не помогает, схема выносит дело в
арбитраж (arbitration), где решение становится обязательным, а
проигравшая сторона платит не только по спорной операции, но и по
процессуальным сборам \sidecite{visa-dispute-rules}\sidecite{mc-chargeback-guide}.

Практический вывод для мерчанта неприятен, но полезен: спор
выигрывает сторона, которая заранее собрала связную доказательную
базу. Логи, чеки, сведения о доставке и согласии клиента должны жить
дольше самого заказа.

\section{Как \scheme{Visa} и \scheme{Mastercard} ведут спор}
\label{sec:disputes-vcr-mc}

В 2018 году \scheme{Visa} перестроила процесс споров на \emph{Visa Claims
Resolution} (VCR)\sidenote{Visa Claims Resolution (VCR) — обновлённая процедура разрешения споров Visa; Order Insight — инструмент раннего снятия спора на уровне банка и карточной схемы. У Mastercard сходную роль играет Consumer Clarity.}.
Смысл реформы состоял в сокращении лишних обменов и в том, чтобы
заставить стороны раньше приносить доказательства
\sidecite{visa-vcr-merchants}\sidecite{visa-dispute-rules}.

Проще всего увидеть эту механику на бытовой покупке. Клиент не узнаёт
списание, звонит в банк, и банк вместо формального чарджбэка сначала
показывает ему дополнительные сведения о покупке. Если мерчант даёт
понятное описание товара, адрес доставки или имя получателя, часть
споров исчезает до формального этапа. Именно на этом и держится
Order Insight: спор пытаются снять ещё до того, как он успеет перейти
в дорогую процессуальную стадию \sidecite{visa-order-insight}.

У \scheme{Mastercard} логика схожа, но язык другой: официальное
руководство говорит о циклах чарджбэка, втором предъявлении
(second presentment) и подаче
дела в арбитраж \sidecite{mc-chargeback-guide}. Иначе говоря,
архитектура очень близка, а словарь свёрстан по-разному.
За прояснение покупки до формального спора в экосистеме
\scheme{Mastercard} сейчас отвечают продукты Ethoca, включая Consumer
Clarity \sidecite{mc-dispute-management}.

Однако сходство не стоит переоценивать. Инструменты раннего снятия
спора помогают только там, где спор возник из-за неузнавания покупки
или недостатка сведений. Если проблема в неполученном товаре, дефектной
услуге или нарушенном условии договора, спор всё равно придётся вести
по полной схеме.

В карточном мире важен не только сам платёж, но и качество его
описания. Чем яснее мерчант обозначен в данных операции, тем меньше у банка
поводов переводить неузнавание покупки в формальный спор.

\section{Коды основания спора: четыре категории \scheme{Visa}}
\label{sec:disputes-reason-codes}

\scheme{Visa} делит коды основания спора на четыре большие категории:
мошенничество (10.x), ошибки авторизации (11.x), ошибки обработки
(12.x) и потребительские споры (13.x) \sidecite{visa-dispute-rules}.
Эта классификация заранее показывает, кто должен доказывать свою
позицию и какие доказательства вообще считаются уместными.

В категории 10.x речь идёт о мошенничестве. Самый важный для интернет-торговли
код здесь~--- 10.4, \textit{Other Fraud -- Card-Absent Environment}:
держатель карты утверждает, что не участвовал в интернет-покупке.
У \scheme{Visa} именно эта категория является главной ареной для
Compelling Evidence 3.0.

В категории 11.x речь идёт о нарушении правил авторизации.
Если мерчант провёл транзакцию без нужного разрешения,
не сохранил корректный ответ авторизации или пошёл дальше, несмотря
на отказ, спор уходит именно сюда.

Категория 12.x описывает ошибки обработки. Здесь спор обычно связан
с некорректной обработкой, а не с намерением клиента: сумма не совпала с
авторизацией, одна и та же операция попала в клиринг дважды, транзакция
попала в схему слишком поздно или с искажёнными данными.

Категория 13.x~--- это потребительские споры: товар не получен, услуга
не оказана, подписка не отменена, возврат обещан, но не проведён.
Здесь важнее всего доказательство исполнения обязательства,
а не технический лог: доставка, акцепт условий, подтверждение пользования,
дата отмены и момент начисления следующего биллингового цикла.

У \scheme{Mastercard} собственная нумерация кодов основания спора, но логика
та же: мошенничество, авторизация, ошибка обработки, спор держателя
карты \sidecite{mc-chargeback-guide}. Разница в номерах не должна
скрывать общую задачу: схема сначала определяет класс спора, а затем
распределяет бремя доказательства.

\section{Убедительные доказательства: что собирать}
\label{sec:disputes-evidence}

Само выражение \enquote{убедительные доказательства} (compelling evidence)
звучит почти бытово, но в споре оно означает очень конкретный набор
фактов, который соответствует конкретному коду основания спора
\sidecite{visa-dispute-rules}.
То, что помогает по мошенничеству, почти никогда не помогает по
неполученному товару.

В споре о дистанционной покупке важно показать привязку транзакции
и к устройству, и к поведению клиента. Сюда относятся IP-адрес,
идентификатор или отпечаток устройства,
учётная запись, адрес доставки, адрес электронной почты, номер телефона,
а при наличии 3-D Secure~--- результат аутентификации. Один IP-адрес
сам по себе слаб; убедительной становится совокупность признаков.

Для спора о неполучении товара или услуги важнее подтверждение
исполнения: отметка курьерской службы о вручении, подпись получателя,
технический лог активации цифрового товара, запись о первом входе в
сервис, протокол оказания услуги. Для спора об отменённой подписке
нужны дата отмены, условия договора и доказательство того, что клиент
видел эти условия до списания.

Документальная гигиена в карточном бизнесе не должна начинаться
после чарджбэка. Логи полезно собирать заранее: IP-адрес, строку
\codeword{User-Agent}, сведения об устройстве, результат AVS и CVV,
результат 3-D Secure, подтверждение доставки и связку между первой
клиентской транзакцией (CIT) и последующими списаниями. Если этой
связки нет, даже сильная по сути позиция в споре быстро слабеет.

\section{Compelling Evidence 3.0: зачем нужен NTID}
\label{sec:disputes-ce30}

С 15 апреля 2023 года \scheme{Visa} расширила правила доказательства, введя
Compelling Evidence 3.0 (CE 3.0)\sidenote{Compelling Evidence 3.0 (CE 3.0) — правило Visa, которое позволяет отбивать спор о мошенничестве в дистанционном канале, если у мерчанта есть две предшествующие неоспоренные транзакции с тем же платёжным реквизитом и совпадающими ключевыми признаками клиента.}.
Это ответ на ситуацию, когда реальный держатель карты оспаривает
собственную покупку так, как будто она была мошеннической
\sidecite{visa-ce30-merchant-readiness}\sidecite{visa-dispute-rules}.

Суть CE 3.0 строже, чем кажется по пересказам. Для защиты нужны две
предыдущие транзакции того же мерчанта без активного сообщения о
мошенничестве (fraud report) и fraud dispute, попадающие в окно
120--365 дней до даты спора; кроме
того, между спорной и историческими транзакциями должны совпасть как
минимум два ключевых признака клиента, и одним из них обязан быть
IP-адрес либо идентификатор устройства
\sidecite{visa-ce30-merchant-readiness}.

Из этого не следует, что \scheme{Visa} требует NTID как отдельный
критерий CE 3.0. Но в подписочном сценарии сохранённый сетевой
идентификатор транзакции (Network Transaction ID, NTID) и аккуратная
связка первой клиентской транзакции (CIT) с последующими списаниями
помогают мерчанту восстановить тот самый исторический след клиента.
Это уже инженерный вывод из сочетания CE 3.0 и Stored Credential
Framework, а не буквальная формулировка правила
\sidecite{visa-ce30-merchant-readiness}\sidecite{visa-scf}.

У \scheme{Mastercard} картина уже не сводится к отсутствию аналога.
Компания развивает First-Party Trust Program: мерчант может передавать
расширенные данные либо в момент авторизации, либо на стадии спора, а
для защиты мерчанта от chargeback программа требует комбинацию
признаков из корзин device factor, delivery factor и additional
identity factors. По терминологии и процессу это не копия
\scheme{Visa} CE~3.0, но и не отсутствие
сопоставимого инструмента
\sidecite{mc-first-party-trust}\sidecite{mc-dispute-management}.

Ограничение CE 3.0 тоже важно помнить. Оно работает только против
споров по мошенничеству в дистанционном канале, а не против неполученного
товара, дефектной услуги или ошибочной отмены подписки. Иными словами,
CE 3.0 помогает там, где спорят о самом факте участия держателя карты,
но не там, где спорят о качестве исполнения обязательства.

\section{Мониторинговые программы схем}
\label{sec:disputes-monitoring}

Системно высокий уровень чарджбэков опасен не только отдельным спором,
но и тем, что мерчант начинает попадать в мониторинговые программы
схем. У этих программ меняются и названия, и пороги, но инвариант один:
если доля споров или мошенничества слишком высока, схема переходит от
предупреждений к плану исправления, штрафам и, в крайнем случае,
к прекращению обслуживания \sidecite{visa-dispute-rules}\sidecite{mc-chargeback-guide}.

В актуальных публичных материалах \scheme{Visa} уже свела мониторинг
споров и мошенничества на уровне эквайера в VAMP, а \scheme{Mastercard}
по-прежнему описывает набор отдельных программ контроля, включая BRAM,
ECP, EFM и QMAP
\sidecite{visa-vamp-intro}\sidecite{mc-compliance-programs}.

Для печатной книги важно понимать механику, а не пытаться запомнить все
актуальные числа. Схема смотрит на сочетание
абсолютного объёма и относительной доли. Сезонность, скользящее окно и
перекрёстная миграция споров между календарными месяцами легко искажают
картину. Поэтому мерчанту полезнее считать свои показатели отдельно по
спорам, мошенничеству и возвратам, чем смотреть на один общий дашборд.

Практический вывод здесь такой: мониторинговая программа не заменяет
доказательства, а лишь подталкивает мерчанта не доводить ситуацию до
массового спора. Чем лучше дескриптор, чище доставка и раньше собраны
доказательства, тем реже схема вообще включает жёсткий режим.

\section{Диспуты в СБП: почему их почти нет}
\label{sec:disputes-sbp}

Держатель карты, недовольный покупкой, нажимает кнопку спора в
банковском приложении, и карточный чарджбэк начинает свой цикл.
Плательщик через \scheme{СБП} такой кнопки обычно не видит не потому,
что система \enquote{недоделана}, а потому, что у неё нет встроенного
карточного механизма принудительного оспаривания уже исполненной
операции.

В СБП деньги зачисляются получателю в режиме реального времени. Если
перевод оказался ошибочным, банк плательщика может запустить процесс
\enquote{Просьба СБП}, но вернуть деньги автоматически нельзя: банки
взаимодействуют с получателем, потому что для возврата требуется его
согласие, если только иной порядок не вытекает из закона, решения суда
или договора с банком \sidecite{cbr-sbp-faq-mistaken-transfer}.
Для оплаты бизнеса при этом существует обычный возврат через банк
продавца, в том числе частичный, и деньги при таком возврате
(refund) поступают
покупателю моментально \sidecite{sbp-faq-business}. Никакого встроенного
сетевого диспута, похожего на карточный, здесь всё равно нет.

Это даёт мерчанту преимущество: ему не грозят классические диспутные
штрафы и длинная цепочка от ответа на чарджбэк до предарбитража.
Для покупателя это означает другое ограничение: если товар не пришёл,
защита строится не на схемном chargeback, а на возврате, который должен
оформить продавец, на переговорах, претензионной работе или обычном
гражданско-правовом споре. Автоматического списания по инициативе банка
покупателя по правилам схемы здесь нет.

Отсюда вывод для продукта. \scheme{СБП}~--- хороший платёжный контур
для быстрого и окончательного перевода, но не замена карточной системе
споров. Поэтому в продукте, который принимает и карты, и СБП, эти два
потока нужно обрабатывать по разным правилам. Единая машина статусов
здесь создаёт лишь ложное ощущение унификации.

\begin{chaptersummary}
\begin{itemize}
  \item Диспут --- это цикл: жалоба, формальный чарджбэк,
        ответ мерчанта, предарбитраж и арбитраж.
  \item \scheme{Visa} и \scheme{Mastercard} описывают похожую механику
        разными словами: \scheme{Visa} говорит о VCR, Order Insight и
        ответе на диспут, а \scheme{Mastercard}~--- о циклах чарджбэка,
        втором предъявлении и подаче дела в арбитраж.
  \item Коды основания спора делят мир на мошенничество, авторизацию,
        ошибки обработки и потребительские споры; один и тот же набор
        доказательств не
        подходит для всех категорий.
  \item \scheme{Visa} CE 3.0 требует исторический след по двум
        предыдущим операциям и совпадающим данным клиента; у
        \scheme{Mastercard} для ложных споров со стороны самого
        держателя карты развивается
        First-Party Trust, но по другой архитектуре.
  \item Схемные мониторинговые программы наказывают не за один спор, а
        за системно высокий уровень споров и мошенничества; важнее не
        цифра сама по себе, а то, как быстро мерчант выстраивает
        собственную защиту.
  \item В СБП нет карточного чарджбэка: ошибочный перевод идёт через
        отдельный процесс банка, а для C2B-оплаты есть обычный refund,
        но не схемный dispute cycle.
\end{itemize}
\end{chaptersummary}
